\section{Related Work}
\label{sec:relwork}

\paragraph{Breach and attack simulation.}
The closest operational analogue to \sysname is continuous adversary
emulation, exemplified by MITRE CALDERA~\cite{caldera} and by a commercial
breach-and-attack-simulation category. These systems run \emph{against
production}, which is their strength, as no fidelity question arises, and
equally their fundamental limitation, as a false step damages a live system
and they are restricted to behaviors known to be safe. They replay
catalogued techniques, and they do not fuzz a parser or exploit an unknown
defect. \sysname makes the opposite trade, accepting a fidelity obligation,
which Section~\ref{sec:thrust2} discharges quantitatively, in exchange for
the freedom to run destructive analysis.

\paragraph{Network configuration verification and attack graphs.}
Batfish~\cite{batfish} and Minesweeper~\cite{minesweeper} build formal models
of network configuration and verify reachability and policy properties over
them, and Batfish in particular is often described as constructing a network
digital twin. This work is the direct ancestor of the logic-based invariants
of Thrust~2, and we adopt its property-checking methodology. The difference
lies in what the model contains, as these tools model the forwarding plane
symbolically, whereas \sysname instantiates \emph{running software} at the
recorded versions and configurations and executes real attacks against it. A
symbolic model can tell us that a host is reachable, but not that reaching it
yields code execution through a defect in the service that runs on it.
Similarly, attack graphs formalize how atomic exploits compose into
breaches~\cite{attackgraphs}, and MulVAL~\cite{mulval} made the approach
practical at deployment scale, but the classical work reasons over a
\emph{database of known} exploits and over assumed host configurations, and the graph is only as good as the
inventory it is given. \sysname
closes that loop by deriving the inventory automatically from observation and
by discovering new primitives through the analysis of the twin.

\paragraph{Discovery, testbeds, and LLM agents for security.}
ZMap~\cite{zmap} and ZGrab~\cite{zgrab} established fast stateless scanning,
osquery~\cite{osquery} exposes host state as queryable tables, and
cartography~\cite{cartography} consolidates infrastructure assets into a
graph database, closely related to the structured layer of Thrust~1. These
tools answer the question of what exists, but none constructs a replica
representative from a security perspective, and none treats the load a probe
imposes on a fragile host as a first-class constraint. Emulation testbeds
such as DETER~\cite{deterlab}, and our own work on GATE~\cite{action}, are
authored by hand, whereas the twins of \sysname are derived automatically
from a specific deployment. Finally, LLM agents have advanced quickly on
offensive tasks~\cite{pentestgpt,cybench,measuring}, but that work evaluates
agents on targets that are \emph{given}, typically single hosts or CTF
challenges, whereas \sysname addresses the problem those agents presuppose,
that is, the construction of the target, at infrastructure scale, and
faithfully enough that what the agent finds is true of a real system.
